% ----------------------------------------------------------------------
% Abstract in English
%
% The abstract provides a concise summary of the research in English.

\selectlanguage{english}

{\centering
\large\bfseries Abstract\par}
% Insere 1,5 linhas de espaço entre o título e o texto do abstract,
% garantindo consistência com o espaçamento utilizado entre títulos
% e texto nas demais seções.
% Insere 1,5cm de espaço vertical entre o título e o texto do abstract.
\vspace{1.5cm}

% O abstract também deve ser digitado em espaço simples conforme as normas
% (citations, footnotes, references, captions etc.).
\begin{singlespace}
Gestational duration ends in two mutually exclusive events, stillbirth and
live birth, which are often analysed separately despite their competing nature.
This dissertation evaluated time to delivery in the state of Sao Paulo,
combining competing risks methods with a Bayesian shared spatial frailty
structure at the level of the Regional Health Care Networks (RRAS). Public
microdata from SIM-DOFET and SINASC for 2023 were harmonized and analysed
after complete-case filtering, resulting in 3,268 stillbirths and 485,922 live
births. Descriptive analyses and cumulative incidence functions showed a clear
asymmetry between outcomes: stillbirths were concentrated at earlier
gestational ages, whereas live births were concentrated near term. Stratified
CIFs indicated relevant social and obstetric patterns, including less favorable
profiles among lower maternal education groups and groups with unfavorable
obstetric history. Indicators by RRAS revealed marked territorial heterogeneity,
with divergence between absolute burden and rates, supporting explicit
modelling of latent regional vulnerability. The proposed Bayesian competing
risks model with ICAR frailty was developed to provide interpretable regional
estimates with uncertainty quantification, supporting maternal and perinatal
health planning within the Brazilian Unified Health System (SUS).

\vspace{0.5cm}
\noindent\textbf{Keywords:} stillbirth. live birth. competing risks. spatial
frailty. Regional Health Care Networks.
\end{singlespace}

% Retorna ao idioma principal (português) após o abstract.
\selectlanguage{brazilian}
